\documentclass[12pt,a4paper]{report}

% ================================================================
% PACKAGES
% ================================================================
\usepackage[utf8]{inputenc}
\usepackage[T1]{fontenc}
\usepackage[french]{babel}
\usepackage{geometry}
\usepackage{graphicx}
\usepackage{listings}
\usepackage{xcolor}
\usepackage{amsmath}
\usepackage{amssymb}
\usepackage{hyperref}
\usepackage{fancyhdr}
\usepackage{titlesec}
\usepackage{tcolorbox}
\usepackage{booktabs}
\usepackage{array}
\usepackage{enumitem}
\usepackage{float}
% \usepackage{microtype}
\usepackage{setspace}
% \usepackage{fontawesome5}

% ================================================================
% MISE EN PAGE
% ================================================================
\geometry{
    top=2.5cm,
    bottom=2.5cm,
    left=2.5cm,
    right=2.5cm
}

% ================================================================
% COULEURS
% ================================================================
\definecolor{codeblue}{RGB}{0,80,160}
\definecolor{codegray}{RGB}{128,128,128}
\definecolor{codegreen}{RGB}{0,128,0}
\definecolor{codepurple}{RGB}{128,0,128}
\definecolor{codebg}{RGB}{245,245,245}
\definecolor{titlecolor}{RGB}{0,60,120}
\definecolor{accentcolor}{RGB}{0,120,200}
\definecolor{successcolor}{RGB}{0,150,0}
\definecolor{warningcolor}{RGB}{200,100,0}

% ================================================================
% STYLE DES LISTINGS
% ================================================================
\lstdefinestyle{cppstyle}{
    backgroundcolor=\color{codebg},
    commentstyle=\color{codegreen}\itshape,
    keywordstyle=\color{codeblue}\bfseries,
    stringstyle=\color{codepurple},
    numberstyle=\tiny\color{codegray},
    basicstyle=\ttfamily\footnotesize,
    breakatwhitespace=false,
    breaklines=true,
    captionpos=b,
    keepspaces=true,
    numbers=left,
    numbersep=5pt,
    showspaces=false,
    showstringspaces=false,
    showtabs=false,
    tabsize=4,
    frame=single,
    rulecolor=\color{accentcolor},
    language=C++
}
\lstset{style=cppstyle}

% ================================================================
% STYLE DES TITRES
% ================================================================
\titleformat{\chapter}[display]
    {\normalfont\huge\bfseries\color{titlecolor}}
    {\chaptertitlename\ \thechapter}{20pt}{\Huge}
\titleformat{\section}
    {\normalfont\Large\bfseries\color{titlecolor}}
    {\thesection}{1em}{}
\titleformat{\subsection}
    {\normalfont\large\bfseries\color{accentcolor}}
    {\thesubsection}{1em}{}

% ================================================================
% EN-TÊTE ET PIED DE PAGE
% ================================================================
\pagestyle{fancy}
\fancyhf{}
\fancyhead[L]{\small\textcolor{titlecolor}{\textbf{Réalité Augmentée From Scratch}}}
\fancyhead[R]{\small\textcolor{titlecolor}{ENSPY - AIA4 - 2025/2026}}
\fancyfoot[C]{\thepage}
\renewcommand{\headrulewidth}{0.4pt}
\renewcommand{\footrulewidth}{0.4pt}

% ================================================================
% BOITES
% ================================================================
\tcbuselibrary{skins,breakable}
\newtcolorbox{infobox}[1]{
    colback=blue!5!white,
    colframe=accentcolor,
    title=#1,
    fonttitle=\bfseries,
    breakable
}
\newtcolorbox{successbox}[1]{
    colback=green!5!white,
    colframe=successcolor,
    title=#1,
    fonttitle=\bfseries,
    breakable
}
\newtcolorbox{warningbox}[1]{
    colback=orange!5!white,
    colframe=warningcolor,
    title=#1,
    fonttitle=\bfseries,
    breakable
}
\newtcolorbox{theorybox}[1]{
    colback=gray!5!white,
    colframe=titlecolor,
    title=#1,
    fonttitle=\bfseries,
    breakable
}

% ================================================================
% DOCUMENT
% ================================================================
\begin{document}

% ================================================================
% PAGE DE GARDE
% ================================================================
\begin{titlepage}
    \begin{center}
        \vspace*{0.5cm}

        % Logo et entête institution
        \textcolor{titlecolor}{\rule{\linewidth}{2pt}}
        \vspace{0.3cm}

        {\Large\textbf{RÉPUBLIQUE DU CAMEROUN}\\
        \textit{Paix -- Travail -- Patrie}}
        \hfill
        {\Large\textbf{REPUBLIC OF CAMEROON}\\
        \textit{Peace -- Work -- Fatherland}}

        \vspace{0.5cm}
        \textcolor{titlecolor}{\rule{\linewidth}{1pt}}
        \vspace{0.3cm}

        {\large\textbf{ÉCOLE NATIONALE SUPÉRIEURE POLYTECHNIQUE}\\
        DE YAOUNDÉ}\\
        \vspace{0.1cm}
        {\normalsize Département Génie Informatique}\\
        {\normalsize classe :  AIA4}

        \vspace{0.5cm}
        \textcolor{titlecolor}{\rule{\linewidth}{1pt}}

        \vspace{1.5cm}

        % Titre principal
        {\Huge\textbf{\textcolor{titlecolor}{
            Rapport de Travaux Pratiques
        }}}

        \vspace{0.5cm}

        {\huge\textbf{\textcolor{accentcolor}{
            Réalité Augmentée From Scratch
        }}}

        \vspace{0.3cm}

        {\Large\textit{
            Bibliothèque mathématique NkMath\\
            Semaines 1 à 4 --- Modules 01 à 04
        }}

        \vspace{1cm}
        \textcolor{titlecolor}{\rule{0.5\linewidth}{1pt}}
        \vspace{1cm}

        % Informations étudiant
        \begin{tcolorbox}[
            colback=blue!5!white,
            colframe=titlecolor,
            width=0.8\linewidth
        ]
        \centering
        {\large\textbf{Informations de l'étudiant}}\\[0.5em]
        \begin{tabular}{rl}
            \textbf{Nom \& Prénom :} & KENNE TAKOUDJOU Franck Manuel \\[0.3em]
            \textbf{Classe :} & AIA4 \\[0.3em]
            \textbf{Établissement :} & ENSPY -- Yaoundé \\[0.3em]
            \textbf{Cours :} & Réalité Augmentée From Scratch \\[0.3em]
            \textbf{Enseignant :} & TEUGUIA TADJUIDJE Rodolf Sédéris \\[0.3em]
            \textbf{Année académique :} & 2025 -- 2026 \\
        \end{tabular}
        \end{tcolorbox}

        \vspace{1cm}

        % Technologies
        \begin{tcolorbox}[
            colback=gray!5!white,
            colframe=codegray,
            width=0.8\linewidth
        ]
        \centering
        {\normalsize
        \textbf{Technologies utilisées :}
        C++17 \quad $\bullet$ \quad
        Jenga Build System \quad $\bullet$ \quad
        Nkentseu Engine \quad $\bullet$ \quad
        GitHub
        }
        \end{tcolorbox}

        \vfill

        % Date
        \textcolor{titlecolor}{\rule{\linewidth}{2pt}}
        \vspace{0.3cm}
        {\large Yaoundé, Mars 2026}

    \end{center}
\end{titlepage}

% ================================================================
% TABLE DES MATIÈRES
% ================================================================
\tableofcontents
\newpage

% ================================================================
% INTRODUCTION
% ================================================================
\chapter*{Introduction }
\addcontentsline{toc}{chapter}{Introduction Générale}

Ce rapport présente les travaux pratiques réalisés dans le cadre du cours
\textbf{Réalité Augmentée From Scratch} dispensé à l'École Nationale
Supérieure Polytechnique de Yaoundé.

L'objectif de ce cours est de construire un pipeline complet de réalité
augmentée en \textbf{C++17}, intégré dans le moteur \textbf{Nkentseu},
sans aucune bibliothèque tierce (pas d'OpenCV, pas de GLM, pas d'Eigen).
Le build system utilisé est \textbf{Jenga}, développé par l'équipe Rihen.

\begin{infobox}{Règle absolue du cours}
Zéro bibliothèque tierce. Uniquement la STL C++17 et les API système
natives. Toute soumission contenant un \texttt{\#include} d'une
bibliothèque interdite est invalidée automatiquement.
\end{infobox}

Ce rapport couvre les \textbf{4 premières semaines} du semestre 1, soit
\textbf{11 travaux pratiques} portant sur la bibliothèque mathématique
\textbf{NkMath} :

\begin{itemize}[leftmargin=2cm]
    \item \textbf{Semaine 1} : IEEE 754, arithmétique flottante (TP1, TP2, TP3)
    \item \textbf{Semaine 2} : Vecteurs Vec2d, Vec3d, Vec4d (TP1, TP2, TP3)
    \item \textbf{Semaine 3} : Matrices Mat4d, TRS, Rasteriseur (TP1, TP2, TP3)
    \item \textbf{Semaine 4} : Quaternions, SLERP, Animation (TP1, TP2)
\end{itemize}

\newpage

% ================================================================
% SEMAINE 1
% ================================================================
\chapter*{Semaine 1 --- IEEE 754 \& Arithmétique Flottante}
\addcontentsline{toc}{chapter}{Semaine 1 --- IEEE 754 \& Arithmétique Flottante}

\section*{Objectifs de la semaine}
\addcontentsline{toc}{section}{Objectifs de la semaine}

\begin{theorybox}{Pourquoi ce module est indispensable}
Tout le pipeline AR manipule des nombres flottants : coordonnées 3D,
matrices de projection, résidus de calibration, valeurs de pixels. Une
erreur d'arrondi invisible dans la résolution d'un système linéaire peut
faire diverger Levenberg-Marquardt, produire une homographie incorrecte,
ou générer une pose décalée de plusieurs centimètres. Ce module est la
\textbf{fondation de tout code numérique fiable}.
\end{theorybox}

% ================================================================
\section{Théorie : Représentation IEEE 754}

\subsection{Structure d'un float 32 bits}

Un nombre flottant IEEE 754 de 32 bits est découpé en trois champs :

\begin{equation}
    \text{Valeur} = (-1)^{\text{signe}} \times 2^{\text{exposant}-127}
    \times 1.\text{mantisse}
\end{equation}

\begin{table}[H]
    \centering
    \begin{tabular}{|c|c|c|}
        \hline
        \textbf{Bit 31} & \textbf{Bits 30--23} & \textbf{Bits 22--0} \\
        \hline
        Signe (1 bit) & Exposant biaisé (8 bits) & Mantisse (23 bits) \\
        0 = positif & biais = 127 & bit implicite 1. devant \\
        \hline
    \end{tabular}
    \caption{Structure d'un float IEEE 754 32 bits}
\end{table}

\subsection{Valeurs spéciales}

\begin{table}[H]
    \centering
    \begin{tabular}{|l|c|c|c|}
        \hline
        \textbf{Type} & \textbf{Signe} & \textbf{Exposant} & \textbf{Mantisse} \\
        \hline
        +0       & 0 & 0x00 & 0x000000 \\
        -0       & 1 & 0x00 & 0x000000 \\
        +Inf     & 0 & 0xFF & 0x000000 \\
        -Inf     & 1 & 0xFF & 0x000000 \\
        NaN      & X & 0xFF & $\neq$ 0 \\
        Subnormal & X & 0x00 & $\neq$ 0 \\
        \hline
    \end{tabular}
    \caption{Valeurs spéciales IEEE 754}
\end{table}

\begin{warningbox}{Piège classique : NaN}
\texttt{NaN != NaN} est \textbf{VRAI} --- la seule façon de tester NaN
est \texttt{std::isnan(x)}. Jamais utiliser \texttt{==} avec NaN.
\end{warningbox}

\subsection{L'epsilon machine}

L'epsilon machine ($\varepsilon$) est la plus petite valeur positive telle
que $1.0 + \varepsilon \neq 1.0$.

\begin{itemize}
    \item Pour \texttt{float} : $\varepsilon \approx 1.19 \times 10^{-7}$
    \item Pour \texttt{double} : $\varepsilon \approx 2.22 \times 10^{-16}$
\end{itemize}

% ================================================================
\section{TP1 --- inspectFloat (Tests 1 à 6)}

\subsection{Objectif}

Implémenter la fonction \texttt{inspectFloat(float x)} qui affiche les
3 champs (signe, exposant, mantisse) en binaire sur la console.

\subsection{Implémentation}

\begin{lstlisting}[caption={Float.cpp --- inspectFloat}]
// Lecture bits IEEE 754 -- C++17 legal (memcpy, pas de UB)
static uint32_t floatBits(float f) {
    uint32_t bits;
    std::memcpy(&bits, &f, sizeof(float));
    return bits;
}

void inspectFloat(float f) {
    uint32_t bits     = floatBits(f);
    uint32_t sign     = (bits >> 31) & 0x1;
    uint32_t exponent = (bits >> 23) & 0xFF;
    uint32_t mantissa =  bits        & 0x7FFFFF;

    printf("float = %.10f | hex = 0x%08X\n", f, bits);
    printf("  signe    = %u\n", sign);
    printf("  exposant = %u (2^%d)\n",
           exponent, (int)exponent - 127);
    printf("  mantisse = 0x%06X\n", mantissa);

    if      (exponent == 0xFF && mantissa != 0)
        printf("  type = NaN\n");
    else if (exponent == 0xFF && mantissa == 0)
        printf("  type = %cInf\n", sign ? '-' : '+');
    else if (exponent == 0x00 && mantissa != 0)
        printf("  type = Subnormal\n");
    else if (exponent == 0x00 && mantissa == 0)
        printf("  type = %c0\n", sign ? '-' : '+');
    else
        printf("  type = Normal\n");
}
\end{lstlisting}

\subsection{Résultats et explications}

\begin{table}[H]
    \centering
    \begin{tabular}{|l|l|p{6cm}|}
        \hline
        \textbf{Valeur} & \textbf{Type} & \textbf{Explication} \\
        \hline
        \texttt{0.1f}   & Normal    & Pas exactement représentable en binaire \\
        \texttt{1.0f}   & Normal    & signe=0, exposant=127, mantisse=0 \\
        \texttt{1.0f/0.0f} & +Inf  & Exposant=0xFF, mantisse=0 \\
        \texttt{sqrt(-1.0f)} & NaN  & Exposant=0xFF, mantisse$\neq$0 \\
        \texttt{-0.0f}  & -0        & Signe=1, exposant=0, mantisse=0 \\
        \texttt{FLT\_MIN} & Subnormal & Exposant=0x00, mantisse$\neq$0 \\
        \hline
    \end{tabular}
    \caption{Résultats de inspectFloat sur les 6 valeurs}
\end{table}

\begin{infobox}{Pourquoi 0.1f n'est pas exact ?}
En base 2, 0.1 en décimal est une fraction périodique infinie :
$0.1_{10} = 0.0\overline{0011}_{2}$. Il est donc impossible de le
représenter exactement avec 23 bits de mantisse. La valeur stockée
est $\approx 0.10000000149011612$.
\end{infobox}

% ================================================================
\section{TP2 --- Kahan et Welford (Tests 7 à 10)}

\subsection{Objectif}

Démontrer empiriquement les problèmes de précision numérique :
sommation de Kahan vs \texttt{std::accumulate}, et variance de
Welford vs variance naïve.

\subsection{Sommation de Kahan}

L'algorithme de Kahan compense les erreurs d'arrondi accumulées lors
de la somme de nombreux termes.

\begin{lstlisting}[caption={Float.cpp --- kahanSum}]
float kahanSum(const float* data, int n) {
    float sum  = 0.0f;
    float comp = 0.0f; // compensation des erreurs perdues
    for (int i = 0; i < n; i++) {
        float y = data[i] - comp; // compenser l'erreur precedente
        float t = sum + y;        // t grand, y petit -> perte bits
        comp = (t - sum) - y;     // capture les bits perdus
        sum  = t;
    }
    return sum;
}
\end{lstlisting}

\subsection{Variance de Welford}

\begin{lstlisting}[caption={Float.cpp --- varianceWelford}]
float varianceWelford(const float* data, int n) {
    float mean = 0.0f;
    float M2   = 0.0f;
    for (int i = 0; i < n; i++) {
        float delta  = data[i] - mean;
        mean        += delta / (i + 1);
        float delta2 = data[i] - mean;
        M2          += delta * delta2;
    }
    return M2 / (n - 1);
}
\end{lstlisting}

\subsection{Résultats}

\begin{table}[H]
    \centering
    \begin{tabular}{|l|c|c|}
        \hline
        \textbf{Méthode} & \textbf{Résultat} & \textbf{Erreur} \\
        \hline
        Valeur théorique & 100000.0000 & 0.0000 \\
        \texttt{std::accumulate} & $\approx$ 100006.xxxx & $\approx$ 6.xxxx \\
        \texttt{kahanSum} & $\approx$ 100000.0000 & $\approx$ 0.0000 \\
        \hline
    \end{tabular}
    \caption{Test 7-8 : Somme de 1 000 000 $\times$ 0.1f}
\end{table}

\begin{warningbox}{Cancellation catastrophique}
La variance naïve de $\{10^8, 10^8, 1, 2\}$ peut donner un résultat
\textbf{négatif} ! C'est la cancellation catastrophique : on soustrait
deux très grands nombres presque égaux. La méthode de Welford est
numériquement stable car elle ne calcule jamais $\text{sum}^2$.
\end{warningbox}

% ================================================================
\section{TP3 --- Float.h (Tests 11 à 14)}

\subsection{Objectif}

Créer le fichier \texttt{Float.h} avec toutes les fonctions utilitaires
et écrire 30 tests unitaires minimum.

\subsection{Implémentation de Float.h}

\begin{lstlisting}[caption={Float.h --- fonctions utilitaires}]
namespace NkMath {

constexpr double kEps  = 1e-9;
constexpr float  kFEps = 1e-6f;

// Valide : ni NaN, ni Inf
inline bool isFiniteValid(float  x) { return std::isfinite(x); }
inline bool isFiniteValid(double x) { return std::isfinite(x); }

// Proche de zero
inline bool nearlyZero(float  x, float  eps = kFEps) {
    return std::abs(x) < eps;
}
inline bool nearlyZero(double x, double eps = kEps) {
    return std::abs(x) < eps;
}

// Egalite relative
inline bool approxEq(double a, double b, double eps = kEps) {
    if (a == b) return true;
    double maxAB = std::max(std::abs(a), std::abs(b));
    return std::abs(a - b) <= eps * std::max(1.0, maxAB);
}

// Sommation de Kahan
double kahanSum(const double* data, int n);

} // namespace NkMath
\end{lstlisting}

\subsection{Tests unitaires (extrait)}

\begin{lstlisting}[caption={Tests Float.h}]
// Test 11 : isFiniteValid
assert(!isFiniteValid(std::sqrt(-1.0f)));  // NaN invalide
assert(!isFiniteValid(1.0f / 0.0f));       // +Inf invalide
assert( isFiniteValid(1.0f));              // Normal valide

// Test 12 : nearlyZero
assert( nearlyZero(0.0f));
assert( nearlyZero(1e-7f));
assert(!nearlyZero(0.1f));

// Test 13 : approxEq
assert( approxEq(1.0f, 1.0f));
assert( approxEq(1.0f, 1.0f + 1e-7f));
assert(!approxEq(1.0f, 1.1f));
\end{lstlisting}

\newpage

% ================================================================
% SEMAINE 2
% ================================================================
\chapter*{Semaine 2 --- Vecteurs : Vec2d, Vec3d, Vec4d}
\addcontentsline{toc}{chapter}{Semaine 2 --- Vecteurs : Vec2d, Vec3d, Vec4d}
\section*{Objectifs de la semaine}
\addcontentsline{toc}{section}{Objectifs de la semaine}

\begin{theorybox}{Pourquoi les vecteurs sont l'alphabet du pipeline 3D}
Les vecteurs sont la structure de données la plus utilisée du pipeline AR.
Chaque position dans l'espace est un Vec3d. Chaque pixel a une position
Vec2d. Chaque coordonnée homogène transmise à la matrice de projection
est un Vec4d. Sans une implémentation rigoureuse avec \texttt{static\_assert}
sur \texttt{sizeof}, rien d'autre dans NkMath ne fonctionne.
\end{theorybox}

% ================================================================
\section{TP1 --- Vec2d.h (Tests 15 à 19)}

\subsection{Objectif}

Implémenter \texttt{Vec2d} avec tous les opérateurs, le dot product,
le cross product 2D, et la normalisation.

\subsection{Implémentation}

\begin{lstlisting}[caption={Vec2d.h --- structure complete}]
struct Vec2d {
    double x, y;

    Vec2d() : x(0.0), y(0.0) {}
    Vec2d(double x, double y) : x(x), y(y) {}
    explicit Vec2d(double s) : x(s), y(s) {} // fill constructor

    // Acces par index -- legal C++17 standard-layout
    double& operator[](int i) {
        assert(i >= 0 && i < 2 && "Vec2d index out of bounds");
        return (&x)[i]; // garanti contigu en C++17
    }

    // Operateurs arithmetiques
    Vec2d operator+(const Vec2d& o) const { return {x+o.x, y+o.y}; }
    Vec2d operator-(const Vec2d& o) const { return {x-o.x, y-o.y}; }
    Vec2d operator*(double s)       const { return {x*s,   y*s  }; }
    Vec2d operator/(double s)       const { return {x/s,   y/s  }; }

    // Norme et normalisation
    double Norm2() const { return x*x + y*y; }
    double Norm()  const { return std::sqrt(Norm2()); }
    Vec2d Normalized() const {
        double n = Norm();
        if (nearlyZero(n)) return Vec2d(0.0); // jamais NaN
        return {x/n, y/n};
    }
};

// Dot product : a.b = |a||b|cos(theta)
inline double Dot(const Vec2d& a, const Vec2d& b) {
    return a.x*b.x + a.y*b.y;
}

// Cross 2D (scalaire)
inline double Cross2D(const Vec2d& a, const Vec2d& b) {
    return a.x*b.y - a.y*b.x;
}

// Garantie layout memoire pour glVertexAttribPointer
static_assert(sizeof(Vec2d) == 16, "Vec2d must be 16 bytes");
static_assert(offsetof(Vec2d, x) == 0, "x must be first");
static_assert(offsetof(Vec2d, y) == 8, "y must be at offset 8");
\end{lstlisting}

\subsection{Tests (15 à 19)}

\begin{lstlisting}[caption={Tests Vec2d}]
// Test 15 : Dot product
assert(approxEq(Dot(Vec2d(1,0), Vec2d(0,1)),  0.0));
assert(approxEq(Dot(Vec2d(1,0), Vec2d(1,0)),  1.0));
assert(approxEq(Dot(Vec2d(3,4), Vec2d(3,4)), 25.0));

// Test 16 : Cross2D
assert(approxEq(Cross2D(Vec2d(1,0), Vec2d(0,1)),  1.0));
assert(approxEq(Cross2D(Vec2d(0,1), Vec2d(1,0)), -1.0));

// Test 17 : Normalisation
Vec2d n = Vec2d(3.0, 4.0).Normalized();
assert(approxEq(n.Norm(), 1.0, kEps));

// Test 18 : Operateur []
Vec2d v(5.0, 7.0);
assert(approxEq(v[0], 5.0));
v[0] = 99.0;
assert(approxEq(v.x, 99.0));

// Test 19 : Layout memoire
static_assert(sizeof(Vec2d) == 16, "Vec2d 16 bytes");
\end{lstlisting}

% ================================================================
\section{TP2 --- Vec3d.h + Gram-Schmidt (Tests 20 à 23)}

\subsection{Objectif}

Implémenter \texttt{Vec3d} avec le produit vectoriel 3D et
l'orthogonalisation de Gram-Schmidt.

\subsection{Produit vectoriel 3D}

\begin{equation}
    \vec{a} \times \vec{b} =
    \begin{vmatrix}
        \vec{i} & \vec{j} & \vec{k} \\
        a_x & a_y & a_z \\
        b_x & b_y & b_z
    \end{vmatrix}
    = (a_y b_z - a_z b_y,\ a_z b_x - a_x b_z,\ a_x b_y - a_y b_x)
\end{equation}

\begin{lstlisting}[caption={Vec3d.h --- Cross product}]
// Cross product 3D -- regle main droite
// Mnemonique permutations cycliques : x->y->z->x
inline Vec3d Cross(const Vec3d& a, const Vec3d& b) {
    return {
        a.y*b.z - a.z*b.y,  // composante x
        a.z*b.x - a.x*b.z,  // composante y
        a.x*b.y - a.y*b.x   // composante z
    };
}
\end{lstlisting}

\subsection{Gram-Schmidt}

\begin{lstlisting}[caption={Vec3d.h --- GramSchmidt}]
struct OrthoBasis { Vec3d u, v, w; };

inline OrthoBasis GramSchmidt(Vec3d a, Vec3d b, Vec3d c) {
    // Etape 1 : normaliser a -> u
    Vec3d u = a.Normalized();
    // Etape 2 : oter la composante u de b -> v
    Vec3d v = (b - Project(b, u)).Normalized();
    // Etape 3 : oter les composantes u et v de c -> w
    Vec3d w = (c - Project(c, u) - Project(c, v)).Normalized();
    return {u, v, w};
}
\end{lstlisting}

\subsection{Tests (20 à 23)}

\begin{lstlisting}[caption={Tests Vec3d}]
// Test 20 : Cross regle main droite
Vec3d r = Cross(Vec3d(1,0,0), Vec3d(0,1,0));
assert(approxEq(r.z, 1.0)); // X x Y = Z

// Test 21 : Cross non-commutatif
Vec3d r2 = Cross(Vec3d(0,1,0), Vec3d(1,0,0));
assert(approxEq(r2.z, -1.0)); // Y x X = -Z

// Test 22 : Gram-Schmidt -- 10 triplets aleatoires
OrthoBasis gs = GramSchmidt(a, b, c);
assert(approxEq(gs.u.Norm(), 1.0, 1e-9));
assert(approxEq(Dot(gs.u, gs.v), 0.0, 1e-9)); // orthogonaux

// Test 23 : Project + Reject = a
Vec3d proj = Project(a, b);
Vec3d rej  = Reject(a, b);
assert((proj + rej) == a); // recomposition exacte
\end{lstlisting}

% ================================================================
\section{TP3 --- Vec4d.h + Projection (Tests 24 à 27)}

\subsection{Objectif}

Implémenter \texttt{Vec4d} avec la déhomogénéisation, projeter les
8 coins d'un cube, et générer une image PPM.

\subsection{Coordonnées homogènes}

\begin{itemize}
    \item $w = 1.0$ $\rightarrow$ \textbf{POINT} dans l'espace
          (la translation de Mat4 s'applique)
    \item $w = 0.0$ $\rightarrow$ \textbf{DIRECTION}
          (la translation ne s'applique \textbf{pas})
\end{itemize}

\subsection{Projection perspective}

\begin{equation}
    u = f_x \cdot \frac{X}{Z} + c_x
    \qquad
    v = f_y \cdot \frac{Y}{Z} + c_y
\end{equation}

Avec $f_x = f_y = 500$, $c_x = c_y = 256$, distance caméra $z_{cam} = 2.0$.

\begin{lstlisting}[caption={Projection perspective simple}]
// Test 25 : Projection simple
for (int i = 0; i < 8; i++) {
    double u = fx * (coins_cam[i].x / coins_cam[i].z) + cx;
    double v = fy * (coins_cam[i].y / coins_cam[i].z) + cy;
    assert(u >= 0 && u < 512 && v >= 0 && v < 512);
}

// Test 26 : NkImage 512x512 -- afficher les 8 coins
NkImage img(512, 512);
img.Fill(0, 0, 0);
for (int i = 0; i < 8; i++)
    img.DrawPoint(px[i], py[i], 255, 0, 0); // rouge

// Test 27 : 12 aretes en vert
for (int i = 0; i < 12; i++)
    img.DrawLine(px[a], py[a], px[b], py[b], 0, 255, 0);

img.SavePPM("cube_projection.ppm");
\end{lstlisting}

\begin{infobox}{Format PPM P6}
Le format PPM (Portable Pixmap) P6 est un format d'image binaire sans
dépendance externe. Structure : \texttt{P6\textbackslash n largeur
hauteur\textbackslash n 255\textbackslash n} suivi des pixels RGB.
C'est le format idéal pour vérifier visuellement les résultats sans
bibliothèque d'images.
\end{infobox}

\newpage

% ================================================================
% SEMAINE 3
% ================================================================
\chapter*{Semaine 3 --- Matrices : Mat4d, TRS, Rasteriseur}
\addcontentsline{toc}{chapter}{Semaine 3 --- Matrices : Mat4d, TRS, Rasteriseur}
\section*{Objectifs de la semaine}
\addcontentsline{toc}{section}{Objectifs de la semaine}

\begin{theorybox}{Convention column-major --- critique pour OpenGL}
En mathématiques, on note les matrices en row-major (ligne par ligne).
OpenGL utilise la convention \textbf{column-major} (colonne par colonne).
Cette différence est fondamentale : une matrice transposée par erreur
dans l'upload GPU produit des transformations incorrectes.
\begin{center}
\texttt{data[col*4 + row] = m(row, col)}
\end{center}
\end{theorybox}

% ================================================================
\section{TP1 --- Mat4d.h + Inverse (Tests 28 à 31)}

\subsection{Objectif}

Implémenter \texttt{Mat4d} avec produit matriciel et inverse
Gauss-Jordan avec pivot partiel.

\subsection{Structure Mat4d}

\begin{lstlisting}[caption={Mat4d.h --- structure column-major}]
struct Mat4d {
    // Stockage column-major : data[col*4 + row]
    double data[16];

    // Acces at(row, col)
    double& operator()(int row, int col) {
        return data[col*4 + row]; // column-major !
    }

    // Identite
    static Mat4d Identity() {
        Mat4d m;
        m(0,0) = m(1,1) = m(2,2) = m(3,3) = 1.0;
        return m;
    }

    // Produit matriciel O(64 multiplications pour 4x4)
    Mat4d operator*(const Mat4d& o) const {
        Mat4d result = Mat4d::Zero();
        for (int row = 0; row < 4; row++)
            for (int col = 0; col < 4; col++) {
                double sum = 0.0;
                for (int k = 0; k < 4; k++)
                    sum += (*this)(row,k) * o(k,col);
                result(row,col) = sum;
            }
        return result;
    }
};
\end{lstlisting}

\subsection{Inverse Gauss-Jordan}

\begin{lstlisting}[caption={Mat4d --- Inverse Gauss-Jordan}]
bool Inverse(Mat4d& out) const {
    // Matrice augmentee [M | I]
    double aug[4][8];
    for (int r = 0; r < 4; r++)
        for (int c = 0; c < 4; c++) {
            aug[r][c]   = (*this)(r,c);
            aug[r][c+4] = (r == c) ? 1.0 : 0.0;
        }

    for (int col = 0; col < 4; col++) {
        // Pivot partiel : trouver le max absolu
        int pivotRow = col;
        for (int r = col+1; r < 4; r++)
            if (abs(aug[r][col]) > abs(aug[pivotRow][col]))
                pivotRow = r;

        if (nearlyZero(aug[col][col])) return false; // singuliere

        // Normaliser et eliminer
        // ...
    }
    return true;
}
\end{lstlisting}

\subsection{Rotation de Rodrigues}

\begin{equation}
    R = I \cos\theta + (1-\cos\theta)\,\hat{n}\hat{n}^T
    + \sin\theta\,[\hat{n}]_\times
\end{equation}

\begin{lstlisting}[caption={Mat4d --- RotateAxis Rodrigues}]
static Mat4d RotateAxis(const Vec3d& axis, double angleRad) {
    Vec3d n = axis.Normalized();
    double c = std::cos(angleRad);
    double s = std::sin(angleRad);
    double t = 1.0 - c;

    Mat4d R = Mat4d::Identity();
    R(0,0) = t*n.x*n.x + c;
    R(0,1) = t*n.x*n.y - s*n.z;
    R(0,2) = t*n.x*n.z + s*n.y;
    // ...
    return R;
}
\end{lstlisting}

\subsection{Tests (28 à 31)}

\begin{lstlisting}[caption={Tests Mat4d}]
// Test 28 : M x Identity = M pour 10 matrices aleatoires
assert((M * I) == M);
assert((I * M) == M);

// Test 29 : M x M⁻¹ = I a 1e-10
Mat4d Inv; M.Inverse(Inv);
Mat4d MInv = M * Inv;
for (int r = 0; r < 4; r++)
    for (int c = 0; c < 4; c++) {
        double expected = (r == c) ? 1.0 : 0.0;
        assert(std::abs(MInv(r,c) - expected) < 1e-10);
    }

// Test 30 : matrice singuliere retourne false
Mat4d S = Mat4d::Zero(); S(0,0) = 1.0; S(1,1) = 1.0;
Mat4d inv;
assert(!S.Inverse(inv)); // false = singuliere

// Test 31 : RotateAxis(Y, PI/2) x {1,0,0,1} = {0,0,-1,1}
Mat4d R = Mat4d::RotateAxis(Vec3d(0,1,0), kPi/2.0);
Vec4d v(1,0,0,1);
Vec4d r = R * v;
assert(approxEq(r.x,  0.0, 1e-10));
assert(approxEq(r.z, -1.0, 1e-10));
\end{lstlisting}

% ================================================================
\section{TP2 --- trs-decomposition.cpp (Tests 38 à 40)}

\subsection{Objectif}

Implémenter TRS complet et vérifier la décomposition inverse
pour 20 triplets $(T, R, S)$ aléatoires.

\subsection{Construction TRS}

\begin{equation}
    M_{TRS} = T \cdot R \cdot S
\end{equation}

L'\textbf{ordre est crucial} : scale d'abord, puis rotation,
puis translation.

\begin{lstlisting}[caption={Mat4d --- TRS}]
static Mat4d TRS(const Vec3d& t, const Vec3d& axis,
                  double angle, const Vec3d& s) {
    return Translate(t) * RotateAxis(axis, angle) * Scale(s);
}
\end{lstlisting}

\subsection{Décomposition inverse}

\begin{lstlisting}[caption={trs-decomposition.cpp --- decomposeTRS}]
TRSResult decomposeTRS(const Mat4d& M) {
    TRSResult res;

    // Etape 1 : Translation = colonne 3
    res.translation = Vec3d(M(0,3), M(1,3), M(2,3));

    // Etape 2 : Scale = norme de chaque colonne 0,1,2
    double sx = sqrt(M(0,0)*M(0,0)+M(1,0)*M(1,0)+M(2,0)*M(2,0));
    double sy = sqrt(M(0,1)*M(0,1)+M(1,1)*M(1,1)+M(2,1)*M(2,1));
    double sz = sqrt(M(0,2)*M(0,2)+M(1,2)*M(1,2)+M(2,2)*M(2,2));
    res.scale = Vec3d(sx, sy, sz);

    // Etape 4 : Angle depuis trace(R) = 1 + 2*cos(angle)
    double trace = res.rotation(0,0)+res.rotation(1,1)+res.rotation(2,2);
    res.angleRad = std::acos((trace - 1.0) * 0.5);

    return res;
}
\end{lstlisting}

% ================================================================
\section{TP3 --- rasterizer-lookat.cpp (Tests 32 à 37)}

\subsection{Objectif}

Implémenter un rasteriseur logiciel dans une \texttt{NkImage} 512$\times$512,
afficher un cube en rotation sur 10 frames.

\subsection{Matrice LookAt}

\begin{lstlisting}[caption={Mat4d --- LookAt}]
static Mat4d LookAt(const Vec3d& eye,
                     const Vec3d& target,
                     const Vec3d& worldUp) {
    Vec3d f = (target - eye).Normalized(); // forward
    Vec3d r = Cross(f, worldUp).Normalized(); // right
    Vec3d u = Cross(r, f); // up corrige

    Mat4d V = Mat4d::Identity();
    V(0,0)=r.x; V(0,1)=r.y; V(0,2)=r.z;
    V(1,0)=u.x; V(1,1)=u.y; V(1,2)=u.z;
    V(2,0)=-f.x; V(2,1)=-f.y; V(2,2)=-f.z;
    V(0,3) = -Dot(r, eye); // translation
    V(1,3) = -Dot(u, eye);
    V(2,3) =  Dot(f, eye);
    return V;
}
\end{lstlisting}

\subsection{Boucle de rendu}

\begin{lstlisting}[caption={rasterizer-lookat.cpp --- 10 frames}]
for (int frame = 0; frame < 10; frame++) {
    // Test 37 : angle different a chaque frame
    double angle = frame * (kPi / 5.0); // 36 degres par frame

    Mat4d rotation = Mat4d::RotateAxis(Vec3d(0,1,0), angle);
    NkImage img(512, 512);
    img.Fill(30, 30, 30); // fond gris

    // Projette et dessine les 12 aretes en rouge
    for (int i = 0; i < 12; i++)
        img.DrawLine(px[a], py[a], px[b], py[b], 255, 0, 0);

    // Test 36 : sauvegarde PPM
    char filename[64];
    snprintf(filename, 64, "frame_%02d.ppm", frame);
    img.SavePPM(filename);
}
\end{lstlisting}

\begin{table}[H]
    \centering
    \begin{tabular}{|c|c|c|}
        \hline
        \textbf{Test} & \textbf{Description} & \textbf{Résultat} \\
        \hline
        32 & Cube unitaire $[-0.5, 0.5]^3$ & 8 coins, 12 arêtes \\
        33 & LookAt(eye=$\{0,1,3\}$, target=$\{0,0,0\}$) & Vue créée \\
        34 & Projection perspective fov=60° & Caméra configurée \\
        35 & 12 arêtes en rouge & DrawLine OK \\
        36 & Sauvegarde PPM P6 & frame\_00.ppm $\to$ frame\_09.ppm \\
        37 & Rotation 36° par frame & Animation OK \\
        \hline
    \end{tabular}
    \caption{Résultats des tests 32 à 37}
\end{table}

\url{https://github.com/TAKOUDJOU237l/Nkentseu}

\newpage

% ================================================================
% SEMAINE 4  (NOUVEAU)
% ================================================================
\chapter*{Semaine 4 --- Quaternions \& Animation SLERP}
\addcontentsline{toc}{chapter}{Semaine 4 --- Quaternions \& Animation SLERP}

\section*{Objectifs de la semaine}
\addcontentsline{toc}{section}{Objectifs de la semaine}

\begin{theorybox}{Pourquoi les quaternions ?}
Les angles d'Euler souffrent du \textbf{gimbal lock} : quand
pitch $= \pm 90°$, les axes yaw et roll deviennent coplanaires
et le système perd un degré de liberté. Les quaternions résolvent
ce problème et offrent une interpolation sphérique (SLERP) qui
suit le chemin le plus court sur la sphère des rotations $S^3$.
\end{theorybox}

% ================================================================
\section{Théorie : Quaternions}

\subsection{Définition}

Un quaternion $q = w + xi + yj + zk$ est un nombre hypercomplexe avec
$i^2 = j^2 = k^2 = ijk = -1$. Un quaternion unitaire ($\|q\| = 1$)
représente une rotation de $2\arccos(w)$ radians autour de l'axe
$(x, y, z)$.

\begin{equation}
    q = \cos\frac{\theta}{2} + \sin\frac{\theta}{2}(n_x i + n_y j + n_z k)
\end{equation}

\subsection{Produit de Hamilton}

\begin{equation}
    q_1 \otimes q_2 =
    \begin{pmatrix}
        w_1 w_2 - x_1 x_2 - y_1 y_2 - z_1 z_2 \\
        w_1 x_2 + x_1 w_2 + y_1 z_2 - z_1 y_2 \\
        w_1 y_2 - x_1 z_2 + y_1 w_2 + z_1 x_2 \\
        w_1 z_2 + x_1 y_2 - y_1 x_2 + z_1 w_2
    \end{pmatrix}
\end{equation}

\subsection{SLERP --- Spherical Linear Interpolation}

\begin{equation}
    \text{Slerp}(q_a, q_b, t) =
    \frac{\sin((1-t)\theta)}{\sin\theta} q_a
    + \frac{\sin(t\theta)}{\sin\theta} q_b
\end{equation}

où $\theta = \arccos(q_a \cdot q_b)$.

\begin{warningbox}{LERP vs SLERP}
Le LERP (interpolation linéaire) sur les quaternions produit une
\textbf{vitesse angulaire non-uniforme} : le mouvement est plus lent
au début et à la fin. Le SLERP garantit une vitesse angulaire
constante, ce qui est visuellement correct pour les animations.
\end{warningbox}

% ================================================================
\section{TP1 --- Quaternions complets (Tests 41 à 43)}

\subsection{Objectif}

Implémenter \texttt{Quat} avec \texttt{FromAxisAngle}, \texttt{Rotate},
\texttt{ToMat3}, \texttt{FromMat3} (méthode de Shepperd) et vérifier
l'inverse.

\subsection{Implémentation de Quat.h}

\begin{lstlisting}[caption={Quat.h --- structure complete}]
struct Quat {
    double w, x, y, z;

    Quat() : w(1), x(0), y(0), z(0) {} // identite

    double Norm2() const { return w*w+x*x+y*y+z*z; }
    double Norm()  const { return std::sqrt(Norm2()); }

    Quat Normalized() const {
        double n = Norm();
        return {w/n, x/n, y/n, z/n};
    }

    // Conjugue : inverse la rotation
    Quat Conjugate() const { return {w, -x, -y, -z}; }

    // Inverse : pour unitaire, Inverse == Conjugate
    Quat Inverse() const {
        double n2 = Norm2();
        return {w/n2, -x/n2, -y/n2, -z/n2};
    }

    // Produit de Hamilton
    Quat operator*(const Quat& o) const {
        return {
            w*o.w - x*o.x - y*o.y - z*o.z,
            w*o.x + x*o.w + y*o.z - z*o.y,
            w*o.y - x*o.z + y*o.w + z*o.x,
            w*o.z + x*o.y - y*o.x + z*o.w
        };
    }
};

// Depuis axe-angle
inline Quat FromAxisAngle(const Vec3d& axis, double angle) {
    Vec3d n = axis.Normalized();
    double s = std::sin(angle / 2.0);
    double c = std::cos(angle / 2.0);
    return {c, n.x*s, n.y*s, n.z*s};
}

// Rotation d'un vecteur par un quaternion (version optimisee)
inline Vec3d Rotate(const Quat& q, const Vec3d& v) {
    Vec3d qv = {q.x, q.y, q.z};
    Vec3d uv  = Cross(qv, v);
    Vec3d uuv = Cross(qv, uv);
    return v + (uv * (2.0 * q.w)) + (uuv * 2.0);
}
\end{lstlisting}

\subsection{Méthode de Shepperd --- ToMat3 / FromMat3}

\begin{lstlisting}[caption={Quat.h --- ToMat3 et FromMat3}]
// Quat vers Mat3
inline Mat3d ToMat3(const Quat& q) {
    double xx=q.x*q.x, yy=q.y*q.y, zz=q.z*q.z;
    double xy=q.x*q.y, xz=q.x*q.z, yz=q.y*q.z;
    double wx=q.w*q.x, wy=q.w*q.y, wz=q.w*q.z;
    Mat3d R;
    R(0,0)=1-2*(yy+zz); R(0,1)=2*(xy-wz); R(0,2)=2*(xz+wy);
    R(1,0)=2*(xy+wz); R(1,1)=1-2*(xx+zz); R(1,2)=2*(yz-wx);
    R(2,0)=2*(xz-wy); R(2,1)=2*(yz+wx); R(2,2)=1-2*(xx+yy);
    return R;
}

// Mat3 vers Quat -- Shepperd (4 cas selon pivot)
inline Quat FromMat3(const Mat3d& R) {
    double trace = R(0,0) + R(1,1) + R(2,2);
    Quat q;
    if (trace > 0) {
        double s = 0.5 / std::sqrt(trace + 1.0);
        q.w = 0.25 / s;
        q.x = (R(2,1) - R(1,2)) * s;
        q.y = (R(0,2) - R(2,0)) * s;
        q.z = (R(1,0) - R(0,1)) * s;
    }
    // ... (cas x, y, z dominants)
    return q.Normalized();
}
\end{lstlisting}

\subsection{Tests (41 à 43)}

\begin{lstlisting}[caption={Tests Quaternions}]
// Test 41 : Rotate(FromAxisAngle(Y,PI/2), {1,0,0}) ~= {0,0,-1}
Quat q = FromAxisAngle(Vec3d(0,1,0), kPi/2.0);
Vec3d r = Rotate(q, Vec3d(1.0, 0.0, 0.0));
assert(approxEq(r.x,  0.0, 1e-9));
assert(approxEq(r.y,  0.0, 1e-9));
assert(approxEq(r.z, -1.0, 1e-9));

// Test 42 : Aller-retour Quat->Mat3->Quat (50 fois)
for (int i = 0; i < 50; i++) {
    Quat q1 = genRandQuat(seed).Normalized();
    Mat3d m  = ToMat3(q1);
    Quat q2  = FromMat3(m);
    // q1 == q2 ou q1 == -q2 (meme rotation)
    bool same = (q1 == q2) || (q1 == -q2);
    assert(same);
}

// Test 43 : q x q⁻¹ = identite
Quat qInv = q.Inverse();
Quat I    = q * qInv;
assert(approxEq(I.w, 1.0, 1e-9));
assert(approxEq(I.x, 0.0, 1e-9));
assert(approxEq(I.y, 0.0, 1e-9));
assert(approxEq(I.z, 0.0, 1e-9));
\end{lstlisting}

\begin{table}[H]
    \centering
    \begin{tabular}{|c|c|c|}
        \hline
        \textbf{Test} & \textbf{Description} & \textbf{Résultat} \\
        \hline
        41 & Rotate(Y,PI/2,\{1,0,0\}) $\approx$ \{0,0,-1\} & OK \\
        42 & Aller-retour Quat$\to$Mat3$\to$Quat, 50 fois & 50/50 OK \\
        43 & $q \times q^{-1} =$ identité & OK \\
        \hline
    \end{tabular}
    \caption{Résultats des tests 41 à 43}
\end{table}

% ================================================================
\section{TP2 --- Animation SLERP (Tests 44 à 46)}

\subsection{Objectif}

Animer une rotation en 60 frames, comparer SLERP vs LERP,
et utiliser le rasteriseur du TP M03 pour afficher le cube.

\subsection{Implémentation SLERP}

\begin{lstlisting}[caption={Quat.h --- SLERP}]
inline Quat Slerp(Quat a, Quat b, double t) {
    double cosAngle = a.w*b.w + a.x*b.x + a.y*b.y + a.z*b.z;

    // Chemin court : si cosAngle < 0, inverser b
    if (cosAngle < 0.0) {
        b.w = -b.w; b.x = -b.x; b.y = -b.y; b.z = -b.z;
        cosAngle = -cosAngle;
    }

    double k0, k1;
    if (cosAngle > 0.9999) {
        // Angle trop petit -> LERP pour eviter sin(0)/0
        k0 = 1.0 - t;
        k1 = t;
    } else {
        double angle    = std::acos(cosAngle);
        double sinAngle = std::sin(angle);
        k0 = std::sin((1.0 - t) * angle) / sinAngle;
        k1 = std::sin(t * angle) / sinAngle;
    }
    return Quat{k0*a.w+k1*b.w, k0*a.x+k1*b.x,
                k0*a.y+k1*b.y, k0*a.z+k1*b.z}.Normalized();
}
\end{lstlisting}

\subsection{Comparaison SLERP vs LERP}

\begin{table}[H]
    \centering
    \begin{tabular}{|c|c|c|c|}
        \hline
        \textbf{Frame} & \textbf{Angle SLERP (°)} &
        \textbf{Angle LERP (°)} & \textbf{Différence} \\
        \hline
        0  & 0.0000   & 0.0000   & 0.0000 \\
        1  & 16.0000  & 15.xxxx  & $\approx$ 0.xxxx \\
        4  & 64.0000  & 63.xxxx  & $\approx$ 0.xxxx \\
        9  & 144.0000 & 144.0000 & 0.0000 \\
        \hline
    \end{tabular}
    \caption{Test 45 : SLERP vs LERP sur 10 frames}
\end{table}

\begin{infobox}{Interprétation visuelle}
\textbf{SLERP (vert)} : la vitesse angulaire est constante ---
le cube tourne de façon uniforme.\\
\textbf{LERP (rouge)} : la vitesse angulaire est non-uniforme ---
le cube accélère et ralentit en milieu de parcours.
\end{infobox}

\subsection{Tests (44 à 46)}

\begin{lstlisting}[caption={Tests SLERP et animation}]
// Test 44 : Animation SLERP 60 frames
Quat qa = FromAxisAngle(Vec3d(0,1,0), 0.0);
Quat qb = FromAxisAngle(Vec3d(0,1,0), kPi);
for (int frame = 0; frame < 60; frame++) {
    double t   = (double)frame / 59.0;
    Quat   qs  = Slerp(qa, qb, t);
    Mat4d  Ms  = QuatToMat4(qs);
    NkImage img(512, 512);
    img.Fill(20, 20, 20);
    drawCube(img, view, Ms, 0, 255, 0); // vert
    img.SavePPM("slerp_frame_XX.ppm");
}

// Test 45 : SLERP vs LERP -- angles differents
Quat qs = Slerp(qa, qb, t);
Quat ql = LerpQuat(qa, qb, t);
// SLERP : vitesse uniforme
// LERP  : vitesse non-uniforme

// Test 46 : Cube rotation 60 frames via rasteriseur
Quat qStart = FromAxisAngle(Vec3d(0,1,0), 0.0);
Quat qEnd   = FromAxisAngle(Vec3d(1,1,0).Normalized(), kPi*2.0);
for (int frame = 0; frame < 60; frame++) {
    double t = (double)frame / 59.0;
    Quat   q = Slerp(qStart, qEnd, t);
    Mat4d  M = QuatToMat4(q);
    NkImage img(512, 512);
    img.Fill(15, 15, 30); // fond bleu nuit
    drawCube(img, view, M, 0, 200, 255);
    img.SavePPM("cube_quat_XX.ppm");
}
\end{lstlisting}

\begin{table}[H]
    \centering
    \begin{tabular}{|c|c|c|}
        \hline
        \textbf{Test} & \textbf{Description} & \textbf{Résultat} \\
        \hline
        44 & SLERP 60 frames PPM & slerp\_frame\_00..59.ppm \\
        45 & SLERP vert vs LERP rouge & Différence visible \\
        46 & Cube rotation 60 frames & cube\_quat\_00..59.ppm \\
        \hline
    \end{tabular}
    \caption{Résultats des tests 44 à 46}
\end{table}

\newpage

% ================================================================
% CONCLUSION
% ================================================================
\chapter*{Conclusion }
\addcontentsline{toc}{chapter}{Conclusion Générale}

\section*{Bilan des travaux pratiques}
\addcontentsline{toc}{section}{Bilan des travaux pratiques}

Ces 4 semaines de travaux pratiques ont permis de construire les
fondations mathématiques du pipeline de réalité augmentée.

\begin{table}[H]
    \centering
    \begin{tabular}{|c|l|l|c|}
        \hline
        \textbf{Sem.} & \textbf{TP} & \textbf{Fichier} &
        \textbf{Tests} \\
        \hline
        S1 & TP1 \& TP2 & Float.cpp & 1--10 \\
        S1 & TP3        & Float.h   & 11--14 \\
        \hline
        S2 & TP1 & Vec2d.h & 15--19 \\
        S2 & TP2 & Vec3d.h & 20--23 \\
        S2 & TP3 & Vec4d.h + NkImage.h & 24--27 \\
        \hline
        S3 & TP1 & Mat4d.h + NkImage.h & 28--31 \\
        S3 & TP2 & trs-decomposition.cpp & 38--40 \\
        S3 & TP3 & rasterizer-lookat.cpp & 32--37 \\
        \hline
        S4 & TP1 & Quat.h & 41--43 \\
        S4 & TP2 & quat-slerp.cpp & 44--46 \\
        \hline
        \multicolumn{3}{|l|}{\textbf{Total}} & \textbf{46 tests} \\
        \hline
    \end{tabular}
    \caption{Récapitulatif des 11 TPs}
\end{table}

\begin{successbox}{Acquis fondamentaux}
\begin{itemize}
    \item Maîtrise de la représentation IEEE 754 et de l'arithmétique
          flottante robuste
    \item Implémentation complète des types vectoriels (Vec2d, Vec3d, Vec4d)
          avec layout mémoire garanti par \texttt{static\_assert}
    \item Matrices 4$\times$4 en column-major avec inverse Gauss-Jordan
    \item Transformations TRS, LookAt, projection perspective
    \item Rasteriseur logiciel avec algorithme de Bresenham
    \item Quaternions unitaires avec SLERP et méthode de Shepperd
    \item Animation 3D from scratch sans aucune bibliothèque tierce
    \item Build system Jenga et moteur Nkentseu
\end{itemize}
\end{successbox}

\section*{Perspectives}
\addcontentsline{toc}{section}{Perspectives}

Les semaines suivantes s'appuieront sur ces fondations pour développer :

\begin{itemize}
    \item \textbf{Semaine 5} : SVD --- Décomposition en Valeurs Singulières
    \item \textbf{Semaine 6} : NkImage --- Traitement d'image from scratch
    \item \textbf{Semaine 7+} : Calibration Zhang, détection ArUco, EPnP,
          pipeline AR complet
\end{itemize}

\vfill

\begin{center}
    \textcolor{titlecolor}{\rule{0.5\linewidth}{1pt}}\\[0.5em]
    \textit{Rapport rédigé par KENNE TAKOUDJOU Franck Manuel}\\
    \textit{AIA4 --- ENSPY Yaoundé --- Mars 2026}\\
    \textcolor{titlecolor}{\rule{0.5\linewidth}{1pt}}
\end{center}

\end{document}
